\ChapterImageStar[cap:justificacion]{Justificación}{./images/fondo.png}\label{cap:justificacion}
\mbox{}\\

\noindentcomponente
En la actualidad, el modelado de Infraestructuras de Tecnologías de la Información (\TI) constituye un componente fundamental en los entornos tecnológicos y organizacionales, debido a que, como afirman~\cite{BROADBENT01}~la misma se reconoce como un factor crítico para la competitividad empresarial. Sin embargo, la discrepancia entre notaciones, metodologías y marcos conceptuales disponibles genera dificultades en la interoperabilidad afectando el intercambio de información y por lo tanto el flujo normal de una organización~\cite{Jinzhi01}. Esta situación repercute directamente en los procesos de enseñanza y aprendizaje del programa de Ingeniería de Sistemas y Computación de la Universidad del Quindío, donde se requiere una aproximación metodológica unificada que permita representar de manera precisa los artefactos de diseño vinculados con la infraestructura de TI.

Tomando en cuenta lo expuesto por~\cite{Božić01} que afirman que una de las ventajas principales del desarrollo basado en modelos es su capacidad de mejorar la comunicación y colaboración entre los equipos de una organizacion al usar un lenguaje común, por lo que se puede intuir que la ausencia del mismo en un entorno academico limita la comunicación efectiva entre docentes, estudiantes e investigadores y dificulta la reutilización y el intercambio de modelos. Adicionalmente~\cite{Peña01}~argumentan que además de los problemas de unificación y heterogeneidad presentados con anterioridad, la complejidad existente en las Infraestructuras de~\TI~hacen aún más necesario el desarrollo de modelos u marcos conceptuales que contribuyan a reducir la dificultad y desalineamiento de las mismas al facilitar la comprensión integral de componentes y dependencias de estas Infraestructuras.

En este contexto, se justifica la propuesta de una notación unificada de modelado para la Infraestructura de TI (NDITI): una notación que optimice la expresividad y la interoperabilidad entre modelos, mejore la coherencia terminológica y potencie prácticas estandarizadas en la docencia y la investigación. La relevancia educativa y formativa del modelado está bien documentada, pues el mismo funge como una herramienta para estudiar y/o verificar la viabilidad, funcionalidad, características, dificultades de diseño y construcción de los sistemas~\cite{Henderson01}, entre ellos los pertenecientes a la infraestructura, por lo que una NDITI adaptada al programa de Ingeniería de Sistemas y Computación contribuiría directamente al fortalecimiento de capacidades académicas e investigativas en la Universidad del Quindío.